\input format.tex

\usepackage{graphicx}
\graphicspath{{cores/}}

\usepackage{environ}
\usepackage{colortbl,array,booktabs}
\usepackage{tabularx}
\usepackage{arydshln}
\newcolumntype{K}[1]{>{\centering\arraybackslash}p{#1}}

\usepackage[nodisplayskipstretch]{setspace}

\colorlet{TablaBordeSuperior}{topcolor}
\colorlet{TablaBordeInferior}{topcolor}
\colorlet{TablaCentroSuperior}{blue!1}
\colorlet{TablaCentroInferior}{blue!20}
\colorlet{FuenteCabeceraTabla}{white}

\newcolumntype{M}[1]{>{\centering\arraybackslash}m{#1}}

\tcbset{rtab/.style={
freelance,
frame code={
 \path[top color=topcolor,bottom color=topcolor]
   ([yshift=-#1*(\baselineskip+2pt)]interior.north west) --
   ([yshift=-#1*(\baselineskip+2pt)]interior.north east) {[sharp corners]--
    ([yshift=3pt]interior.north east) --
    ([yshift=3pt]interior.north west)} -- cycle;

  },
interior code={},
 }
}

\newcommand\fuentecabecera[1]{\textcolor{black}{\textbf{#1}}}


\begin{document}

\vspace*{4mm}
%% 各章节
\color{black70}
\noindent\fontsize{16pt}{18pt}\selectfont{\color{topcolor} 检测结果}

\vspace*{3mm}
\begin{spacing}{1.5}
\noindent{\fontsize{10pt}{11pt}\selectfont\bf {1、主要益生菌检测结果}}
\end{spacing}

\arrayrulecolor{gray2}
\setlength{\arrayrulewidth}{.5pt}
\fontsize{9pt}{11pt}\selectfont

\smallskip
%{
\lantxh
\begin{tctabularx}{tabularx={m{4.8cm}<{\centering}m{5.2cm}<{\centering}m{4.55cm}<{\centering}}}
&&
\\[-6pt]
\cellcolor{topcolor} \raisebox{2.614pt}{\color{white}\bf 主要益生菌属} &
\cellcolor{topcolor} \raisebox{2.614pt}{\color{white}\bf 本次检测值} &
\cellcolor{topcolor} \raisebox{2.614pt}{\color{white}\bf 参考范围}
\end{tctabularx}

\vspace*{-4.25mm}
\fontsize{8.8pt}{11pt}\selectfont
\begin{longtable}{|m{4.8cm}<{\centering}|m{5.2cm}<{\centering}|m{4.53cm}<{\centering}|}
\hline
\parbox[c]{\hsize}{\vskip7pt 乳酸杆菌属（Lactobacillus） \vskip7pt} & 0.01 
{\includegraphics[scale=1]{redarrow_down.pdf}}
& ≥ 0.04 \\
\hline
\parbox[c]{\hsize}{\vskip7pt 双歧杆菌属（Bifidobacterium） \vskip7pt} & 5.8 
& ≥ 0.21 \\
\hline
\parbox[c]{\hsize}{\vskip7pt 阿克曼氏菌属（Akkermansia） \vskip7pt} & 0 
{\includegraphics[scale=1]{redarrow_down.pdf}}
& ≥ 0.01 \\
\hline
\parbox[c]{\hsize}{\vskip7pt 柔嫩梭菌属（Faecalibacterium） \vskip7pt} & 204.63 
& ≥ 27.24 \\
\hline
\end{longtable}

\vspace*{-2.5mm}
\noindent {结果说明：\includegraphics[scale=1]{redarrow_down.pdf} \ 检测值相比参考人群偏低}

\vspace*{3mm}
\begin{spacing}{1.5}
\noindent{\fontsize{10pt}{11pt}\selectfont\bf {2、检测结果分析}}

\vspace*{2mm}
{\indent\fontsize{9pt}{10pt}\selectfont 
您本次检测的4种主要益生菌含量相比参考人群均偏低。乳酸杆菌作为最知名的益生菌，是调节肠道环境的重要细菌，同时增强免疫力、抵抗病原菌、促进营养吸收等都有它的功劳。双歧杆菌作为肠道内的“清道夫”，能调节肠道微生态平衡、增强肠道屏障功能、抑制有害菌增殖。阿克曼氏菌是能够帮助控制体重的“瘦菌”，它促进脂肪代谢，消耗多余能量，维持代谢平衡。柔嫩梭菌是健康成人肠道中含量较高的常见有益菌，它能产生大量有益物质丁酸，帮助保护肠道、抑制炎症。
这些菌属含量偏低，有可能导致免疫调节力下降、肠道环境差、抵抗病原菌能力降低等，
甚至增加胃肠道疾病、炎症性肠病、腹痛腹泻的风险。
\\\\
\indent 需要注意的是，肠道菌群处于动态平衡，单一菌种数量的短期改变并非致病的直接原因，益生菌的长期异常，整体肠道菌群失衡，才可能引起疾病的发生。}
\end{spacing}

\vspace*{5mm}

\noindent\colorbox{topcolor}{
\parbox{0.6\hsize}{
\color{white}\fontsize{9pt}{\baselineskip}\selectfont { 以上为您本次检测结果，如有任何疑问，请咨询专业医生！}
}
}

\vspace*{.5mm}

\noindent\colorbox{topcolor}{
\parbox{0.6\hsize}{
\color{white}\fontsize{9pt}{\baselineskip}\selectfont { 再次感谢您的信任和支持！}
}
}

\vspace*{3mm}

\hfill\fontsize{9pt}{11pt}\selectfont 人和未来生物科技（长沙）有限公司

\vspace*{-2mm}

\hfill\includegraphics[scale=0.6]{erweima.pdf}


\hfill\fontsize{9pt}{11pt}\selectfont 扫码关注“常道康”

\end{document}
